\section{Evaluation}

\subsection{LLM-as-judge prompts}
\label{appendix:llm_judge_prompt}
\subsubsection{Visual Tasks}
We used GPT5.2~\cite{gpt-5.2} for visual tasks labeling. We give the LLM judge the chunked source as well as chunked model generation and ask it to treat it as a multi-label prediction problem.

\begin{lstlisting}
prompt = '''You are an annotation assistant for SOURCE ATTRIBUTION.

You will be given:
1) An image that contains the source conversation with image IDs shown on the images.
2) A model-generated summary (split into sentences, each sentence has a sentence index).
3) The source text chunks listed below, each with a numeric text ID.

Task:
For EACH generated summary sentence, decide which source evidence supports it.
- "text_source": a list of TEXT IDs from the provided text chunks that directly support that sentence.
- "image_source": a list of IMAGE IDs visible in the provided image that directly support that sentence.
Return ONLY IDs that are explicitly supported by the source. Do NOT invent IDs.

Attribution rules:
- Include an ID only if it is necessary to justify the sentence content.
- If the sentence is supported by multiple text or image chunks, include all relevant IDs.
- If the sentence refers to visual content (objects, actions, clothing, scene, etc.), include the relevant image IDs.
- If a sentence is NOT directly supported by any source evidence, output an empty list for that field.
- Do NOT use the instruction chunk (e.g., "Summarize the conversation") as evidence unless the sentence is literally about the instruction.
- Do NOT merge evidence across sentences: each sentence gets its own dict.
- IDs must be integers (no quotes), in ascending order, with no duplicates.

Output format (STRICT):
Return a JSON array with length == number of generated summary sentences.
Each element corresponds to the sentence index order (0..N-1) and must be exactly:
{{"image_source": [...], "text_source": [...]}}

Do not output any extra keys, explanations, markdown, or trailing text.

Model generated response (sentences indexed):
{model_generation}

Text source chunks (format: [text_id] text):
{text_source_chunks}
'''
\end{lstlisting}

\subsubsection{Audio and Video Tasks}
Since Audio and Video tasks has continuous labels in timeline, we are not able to provide chunked source as in the visual-text setting. We give the judge the raw audio/video source and ask it output relevant timespans for each generated chunk.

\begin{lstlisting}
prompt = '''You are an annotation assistant for SOURCE ATTRIBUTION.

You will be given:
1) One AUDIO/VIDEO file that contains the source content.
2) A model-generated summary (split into sentences, each sentence has a sentence index).

Task:
For EACH generated summary sentence, decide which source evidence supports it.
- "audio/video_source": a list of timestamp spans [start_sec, end_sec] in the AUDIO/VIDEO that directly support that sentence.

Timestamp labeling requirements:
- Timestamps are continuous spans in SECONDS.
- Use a list of spans to allow multiple supporting regions, e.g. [[0,2], [6,8]].
- Each span must satisfy: start_sec < end_sec.
- Use minimal coverage: include ONLY the smallest time ranges that contain the relevant evidence.
  - Do NOT label the entire audio (e.g., [0,100]) unless the sentence truly requires the entire audio.
- Prefer a small number of short spans rather than one very long span.
- Round timestamps to 1s precision and keep them consistent.

Attribution rules:
- Include a timestamp span only if it is necessary to justify the sentence content.
- If the sentence is supported by multiple audio/video regions, include all relevant evidence.
- If a sentence is NOT directly supported by any source evidence in a modality, output an empty list for that field.
- Do NOT use the instruction chunk (e.g., "Summarize the conversation") as evidence unless the sentence is literally about the instruction.
- Do NOT merge evidence across sentences: each sentence gets its own dict.
- audio/video_source spans must be sorted by start time, non-overlapping when possible (merge overlaps if they touch or overlap).

Output format (STRICT):
Return a JSON array with length == number of generated summary sentences.
Each element corresponds to the sentence index order (0..N-1) and must be exactly:
{{"audio/video_source": [[start,end], ...]}}

Do not output any extra keys, explanations, markdown, or trailing text.

Model generated response (sentences indexed):
{model_generation}

STRICT OUTPUT CONTRACT:
- Output MUST be valid JSON only (no markdown, no code fences).
- Do not include ```json or ``` anywhere.
- Output must start with '[' and end with ']'.
- No extra text before or after the JSON.
- Keep it compact on one line.
'''
\end{lstlisting}

\subsection{Human study}
\label{appendix:human_study}
To verify the quality of LLM judge labels, we hired four human experts to manually go through around 26.6\% of the data using Qwen2.5-Omni responses, and report the agreement using the same metric: span-level F1 for visual and time F1 for audio and video. Results are reported in \Cref{tab:human_llm_agreement}.

\begin{table}[ht]
\centering
\tablesize
\caption{Agreement between human annotations and LLM-judge labels on a subset of the evaluation data (26.6\%). Metrics follow the same definitions used in the main evaluation: span-level F1 for visual attribution and time-based F1 for audio/video attribution.}
\label{tab:human_llm_agreement}
\skiptablebegin
\begin{tabular}{lccccc}
\toprule
\textbf{Agreement} & \multicolumn{3}{c}{\textbf{Visual}} & \textbf{Audio} & \textbf{Video} \\
\cmidrule(lr){2-4} \cmidrule(lr){5-5} \cmidrule(lr){6-6}
\textbf{Task} & \multicolumn{2}{c}{Summarization} & QA & QA & QA \\
\textbf{Metric} & Text F1 & Image F1 & Image F1 & Time F1 & Time F1 \\
\midrule
Score & 83.04 & 94.30 & 93.80 & 81.75 & 85.28 \\
\# Samples & 61 & 61 & 40 & 40 & 40 \\
\bottomrule
\end{tabular}
\skiptableend
\end{table}

\subsection{Faithfulness Verification}
\label{appendix:faithfulness_verification}
Beyond attribution accuracy with respect to external labels, we evaluate whether OmniTrace reflects the model’s own decision process.
Specifically, for multiple-choice QA tasks, we test whether the attribution assigned to the model’s final answer is consistent with the option it selects.

\paragraph{Option-consistency evaluation.}
All QA datasets are formulated as multiple-choice questions, where each option corresponds to a discrete source text chunk.
For each sample, we extract the model’s predicted option $\hat{o}$ from the generated response using deterministic parsing. Importantly, this evaluation does not depend on whether the prediction is correct; it tests whether attribution reflects the model’s chosen answer, not whether the answer matches ground truth.

Let $C_k$ denote the generation span containing the final answer statement. For each option $o$, let $S_o$ denote the source chunk corresponding to that option. Using the token-level mappings and confidence-weighted votes defined in \Cref{sec:problem}, we compute the option-restricted attribution mass:
\[
\text{Mass}(S_o)
=
\sum_{t \in T_k}
w_t \cdot \mathbf{1}[\hat{s}(t)=S_o],
\]
where $T_k$ are tokens in the answer span, $\hat{s}(t)$ is the source traced for token $t$,
and $w_t$ is the confidence- and POS-weighted vote. We then define the predicted attribution option as
\[
\tilde{o}
=
\arg\max_{o} \text{Mass}(S_o).
\]

\paragraph{Top-1 Option Consistency.}
We report \emph{Top-1 consistency}, defined as
\[
\text{Consistency}
=
\mathbb{1}[\tilde{o} = \hat{o}],
\]
averaged over samples.
This metric measures whether the option receiving the highest attribution mass matches the option explicitly selected by the model.

\paragraph{Results.}
On the 200 visual QA samples, OmniTrace (Qwen + OT$_\text{AttMean}$) achieves \textbf{93.84\%} Top-1 consistency. This indicates that in the vast majority of cases, the decision sentence’s attribution mass is concentrated on the source chunk corresponding to the option chosen by the model.
In other words, the generation-time tracing mechanism recovers the model’s own selected option as the most supported source unit.

\paragraph{Interpretation.}
This result provides strong evidence that OmniTrace captures attribution signals aligned with the model’s internal decision pathway.
If attribution were dominated by post-hoc noise or unrelated spans, the consistency rate would approach chance level (25\% for four-way multiple choice).
Instead, the observed 93.84\% rate suggests that generation-time token tracing, combined with span-level aggregation and filtering, faithfully identifies the textual source corresponding to the model’s final answer statement.
